\section{应记应背}

\subsection{分子与细胞}

\subsubsection{组成细胞的分子}

\begin{itemize}
  \item \textbf{大量元素}：C、H、O、N、P、S、K、Ca、Mg
  \item \textbf{微量元素}：Fe、Mn、B、Zn、Mo、Cu（铁锰硼锌钼铜）
  \item \textbf{最基本元素}：C（因碳链构成生物大分子骨架）
  \item \textbf{水的存在形式}：自由水（良好溶剂、参与反应、运输）和结合水（细胞结构成分）
  \item \textbf{无机盐作用}：构成化合物（Fe$^{2+}$--血红蛋白、Mg$^{2+}$--叶绿素）、维持渗透压和酸碱平衡
  \item \textbf{糖类}：单糖（葡萄糖、果糖、半乳糖、核糖、脱氧核糖）、二糖（蔗糖、麦芽糖、乳糖）、多糖（淀粉、纤维素、糖原、几丁质）
  \item \textbf{脂质}：脂肪（储能、保温、缓冲）、磷脂（生物膜主要成分）、固醇（胆固醇、性激素、维生素D）
  \item \textbf{蛋白质}：基本单位氨基酸（20种，8种必需氨基酸），脱水缩合形成肽键（--CO--NH--）
  \begin{itemize}
    \item 结构层次：氨基酸 $\to$ 多肽 $\to$ 盘曲折叠 $\to$ 蛋白质
    \item 功能：结构蛋白（羽毛、肌肉）、催化（酶）、运输（血红蛋白）、信息传递（胰岛素）、防御（抗体）
  \end{itemize}
  \item \textbf{核酸}：DNA 和 RNA，基本单位核苷酸（一分子磷酸+五碳糖+含氮碱基）
  \begin{itemize}
    \item DNA 特有：脱氧核糖、T（胸腺嘧啶）
    \item RNA 特有：核糖、U（尿嘧啶）
    \item 核酸功能：携带遗传信息，控制蛋白质合成
  \end{itemize}
  \item \textbf{生物大分子}：多糖、蛋白质、核酸——以碳链为骨架
\end{itemize}

\subsubsection{细胞的结构与功能}

\begin{itemize}
  \item \textbf{原核细胞 vs 真核细胞}
  \begin{itemize}
    \item 原核：无核膜包被的细胞核，只有核糖体（唯一细胞器），如细菌、蓝藻
    \item 真核：有核膜包被的细胞核，有多种细胞器
  \end{itemize}
  \item \textbf{细胞膜}：流动镶嵌模型，磷脂双分子层（基本骨架）+ 蛋白质（嵌入或贯穿）+ 糖蛋白（识别）
  \begin{itemize}
    \item 功能：将细胞与外界分隔、控制物质进出、进行细胞间信息交流
  \end{itemize}
  \item \textbf{细胞器}
  \begin{itemize}
    \item 线粒体：有氧呼吸主要场所，含少量 DNA，双层膜
    \item 叶绿体：光合作用场所，含 DNA 和色素，双层膜
    \item 内质网：蛋白质合成加工、脂质合成，单层膜
    \item 高尔基体：加工、分类、包装蛋白质（植物细胞中与细胞壁形成有关），单层膜
    \item 核糖体：蛋白质合成场所，无膜结构（rRNA+蛋白质）
    \item 溶酶体：含水解酶，消化分解，单层膜
    \item 液泡：调节细胞渗透压，单层膜
    \item 中心体：与有丝分裂有关（动物细胞、低等植物），无膜结构
  \end{itemize}
  \item \textbf{细胞核}：核膜（双层膜）、核孔（大分子通道）、核仁（rRNA 合成）、染色质（DNA+蛋白质）
\end{itemize}

\subsubsection{物质跨膜运输}

\begin{itemize}
  \item \textbf{被动运输}：顺浓度梯度，不耗能
  \begin{itemize}
    \item 自由扩散：O$_2$、CO$_2$、H$_2$O、甘油、乙醇、苯
    \item 协助扩散：葡萄糖（进入红细胞）、离子（通道蛋白），需载体/通道蛋白
  \end{itemize}
  \item \textbf{主动运输}：逆浓度梯度，需载体蛋白，耗能（ATP），如 K$^+$、Na$^+$、葡萄糖（进入小肠上皮细胞）
  \item \textbf{胞吞胞吐}：大分子颗粒，依赖膜的流动性，耗能
\end{itemize}

\subsubsection{酶与 ATP}

\begin{itemize}
  \item \textbf{酶}：多数是蛋白质（少数是 RNA），高效性、专一性、作用条件温和
  \item 酶催化机理：降低反应活化能
  \item 影响酶活性因素：温度、pH、底物浓度、酶浓度
  \item \textbf{ATP}：腺苷三磷酸（A--P$\sim$P$\sim$P），直接能源物质
  \item ATP $\rightleftharpoons$ ADP + Pi + 能量（ATP 合成酶催化合成，ATP 水解酶催化水解）
\end{itemize}

\subsubsection{细胞呼吸}

\begin{itemize}
  \item \textbf{有氧呼吸}（主要）：C$_6$H$_{12}$O$_6$ + 6O$_2$ $\to$ 6CO$_2$ + 6H$_2$O + 能量
  \begin{itemize}
    \item 第一阶段（细胞质基质）：葡萄糖 $\to$ 丙酮酸 + [H] + 少量 ATP
    \item 第二阶段（线粒体基质）：丙酮酸 + H$_2$O $\to$ CO$_2$ + [H] + 少量 ATP
    \item 第三阶段（线粒体内膜）：[H] + O$_2$ $\to$ H$_2$O + 大量 ATP
  \end{itemize}
  \item \textbf{无氧呼吸}
  \begin{itemize}
    \item 乳酸发酵：葡萄糖 $\to$ 2 乳酸 + 少量能量（动物、马铃薯块茎）
    \item 酒精发酵：葡萄糖 $\to$ 2 酒精 + 2CO$_2$ + 少量能量（植物、酵母菌）
  \end{itemize}
  \item 细胞呼吸意义：为生命活动供能，为物质转化提供中间产物
\end{itemize}

\subsubsection{光合作用}

\begin{itemize}
  \item 总反应：6CO$_2$ + 12H$_2$O $\xrightarrow{\text{光能、叶绿体}}$ C$_6$H$_{12}$O$_6$ + 6O$_2$ + 6H$_2$O
  \item \textbf{光反应阶段}（类囊体薄膜）
  \begin{itemize}
    \item 水的光解：H$_2$O $\to$ [H] + O$_2$
    \item ATP 合成：ADP + Pi $\xrightarrow{\text{光能}}$ ATP
    \item 条件：光、色素、酶
  \end{itemize}
  \item \textbf{暗反应阶段}（叶绿体基质）
  \begin{itemize}
    \item CO$_2$ 固定：CO$_2$ + C$_5$ $\to$ 2C$_3$
    \item C$_3$ 还原：2C$_3$ + [H] + ATP $\to$ (CH$_2$O) + C$_5$
    \item 条件：酶、ATP、[H]（不直接需光）
  \end{itemize}
  \item 影响光合作用因素：光照强度、CO$_2$ 浓度、温度、水分、矿质元素
  \item 化能合成作用：利用无机物氧化释放的化学能将 CO$_2$ 和 H$_2$O 合成有机物（如硝化细菌）
\end{itemize}

\subsubsection{细胞增殖}

\begin{itemize}
  \item \textbf{有丝分裂}（体细胞增殖方式）
  \begin{itemize}
    \item 间期：DNA 复制 + 蛋白质合成（中心体复制）
    \item 前期：染色质 $\to$ 染色体，核膜核仁消失，纺锤体形成
    \item 中期：染色体排列于赤道板，数目形态最清晰（观察最佳时期）
    \item 后期：着丝粒分裂，姐妹染色单体分开移向两极
    \item 末期：染色体 $\to$ 染色质，核膜核仁重现，纺锤体消失，细胞质分裂
  \end{itemize}
  \item \textbf{减数分裂}（生殖细胞形成）
  \begin{itemize}
    \item 间期：DNA 复制
    \item 减 I：同源染色体联会（四分体）、交叉互换、同源染色体分离、非同源染色体自由组合
    \item 减 II：着丝粒分裂，姐妹染色单体分离（类似有丝分裂）
  \end{itemize}
  \item 有丝分裂 vs 减数分裂：有丝分裂产生体细胞（2n$\to$2n），减数分裂产生配子（2n$\to$n）
  \item \textbf{细胞分化}：基因选择性表达的结果，细胞在形态、结构和功能上发生稳定性差异
  \item \textbf{细胞全能性}：已分化的细胞仍具有发育成完整个体的潜能
  \item \textbf{细胞衰老}：水分减少、酶活性降低、色素积累、呼吸减慢、核增大
  \item \textbf{细胞凋亡}：基因程序性死亡（主动），对生物体有利
  \item \textbf{细胞坏死}：外界不利因素导致（被动）
\end{itemize}

\subsection{遗传与进化}

\subsubsection{遗传的细胞基础}

\begin{itemize}
  \item 精子的形成：精原细胞 $\to$ 初级精母细胞 $\to$ 次级精母细胞 $\to$ 精细胞 $\to$ 精子
  \item 卵细胞形成：卵原细胞 $\to$ 初级卵母细胞 $\to$ 次级卵母细胞 + 第一极体 $\to$ 卵细胞 + 第二极体
  \item 受精作用：精核与卵核融合，染色体数目恢复为 2n
\end{itemize}

\subsubsection{遗传的分子基础}

\begin{itemize}
  \item \textbf{DNA 是遗传物质的证据}
  \begin{itemize}
    \item 肺炎链球菌转化实验（格里菲斯：转化因子；艾弗里：DNA 是转化因子）
    \item 噬菌体侵染细菌实验（赫尔希和蔡斯，$^{35}$S 标记蛋白质，$^{32}$P 标记 DNA）
  \end{itemize}
  \item \textbf{DNA 结构}
  \begin{itemize}
    \item 双螺旋结构（沃森和克里克）
    \item 外侧：磷酸-脱氧核糖交替排列
    \item 内侧：碱基互补配对（A--T 双键，C--G 三键）
    \item 方向：反向平行
  \end{itemize}
  \item \textbf{DNA 复制}：半保留复制，需解旋酶、DNA 聚合酶、原料（dNTP）、模板
  \item \textbf{转录}：DNA $\to$ RNA，RNA 聚合酶催化，场所细胞核
  \item \textbf{翻译}：mRNA $\to$ 蛋白质，核糖体、tRNA（反密码子识别密码子）
  \item \textbf{中心法则}
  \begin{itemize}
    \item DNA $\xrightarrow{\text{复制}}$ DNA
    \item DNA $\xrightarrow{\text{转录}}$ RNA $\xrightarrow{\text{翻译}}$ 蛋白质
    \item RNA $\xrightarrow{\text{复制}}$ RNA（某些病毒）
    \item RNA $\xrightarrow{\text{逆转录}}$ DNA（HIV 等逆转录病毒）
  \end{itemize}
  \item \textbf{基因对性状的控制}
  \begin{itemize}
    \item 直接控制：基因通过控制蛋白质结构直接控制性状
    \item 间接控制：基因通过控制酶的合成控制代谢过程进而控制性状
  \end{itemize}
\end{itemize}

\subsubsection{遗传规律}

\begin{itemize}
  \item \textbf{分离定律}：等位基因在减 I 后期随同源染色体分离而分开
  \item \textbf{自由组合定律}：非同源染色体上的非等位基因自由组合
  \item \textbf{常见比例}
  \begin{itemize}
    \item 分离定律 3:1（杂合子自交），1:1（测交）
    \item 自由组合 9:3:3:1（双杂自交），1:1:1:1（双杂测交）
    \item 9:7（互补作用），9:3:4（隐性上位），12:3:1（显性上位），15:1（重叠作用），9:6:1（累加作用）
  \end{itemize}
  \item \textbf{伴性遗传}
  \begin{itemize}
    \item 伴 X 显性：女性患者多于男性，可交叉遗传（抗维生素 D 佝偻病）
    \item 伴 X 隐性：男性患者多于女性，交叉遗传（色盲、血友病）
    \item 伴 Y 遗传：限雄遗传（外耳道多毛症）
  \end{itemize}
  \item \textbf{判断遗传方式口诀}：无中生有为隐性，隐性遗传看女病，父子皆病为伴性；有中生无为显性，显性遗传看男病，母女皆病为伴性
\end{itemize}

\subsubsection{变异与育种}

\begin{itemize}
  \item \textbf{基因突变}：DNA 碱基对的替换、增添、缺失
  \begin{itemize}
    \item 特点：普遍性、随机性、不定向性、低频性、多害少利性
    \item 诱变因素：物理（射线）、化学（亚硝酸）、生物（病毒）
  \end{itemize}
  \item \textbf{基因重组}：减 I 前期的交叉互换、减 I 后期的自由组合、转基因
  \item \textbf{染色体变异}
  \begin{itemize}
    \item 结构变异：缺失、重复、倒位、易位
    \item 数目变异：个别增减（单体、三体）、成倍增减（多倍体）
    \item 染色体组：细胞中形态功能各不相同的一组非同源染色体
  \end{itemize}
  \item \textbf{育种方法}
  \begin{itemize}
    \item 杂交育种：基因重组，操作简单，周期长
    \item 诱变育种：基因突变，提高突变率，有利变异少
    \item 单倍体育种：花药离体培养 + 秋水仙素处理，缩短育种年限
    \item 多倍体育种：秋水仙素处理萌发的种子/幼苗，器官大，结实率低
    \item 基因工程育种：定向改造生物性状
  \end{itemize}
  \item \textbf{单倍体}：由配子直接发育而来，体细胞含本物种配子染色体数目
  \item \textbf{多倍体}：体细胞含三个或以上染色体组
\end{itemize}

\subsubsection{进化}

\begin{itemize}
  \item \textbf{现代生物进化理论}
  \begin{itemize}
    \item 种群是生物进化的基本单位
    \item 进化的实质：种群基因频率的改变
    \item 突变和基因重组为进化提供原材料
    \item 自然选择决定进化的方向
    \item 隔离是物种形成的必要条件
  \end{itemize}
  \item \textbf{基因频率计算}：$p+q=1$，$p^2+2pq+q^2=1$（Hardy--Weinberg 平衡）
  \item \textbf{物种形成方式}：地理隔离 $\to$ 生殖隔离
  \item \textbf{共同进化}：不同物种之间、生物与无机环境之间相互影响中不断进化发展
\end{itemize}

\subsection{稳态与调节}

\subsubsection{内环境与稳态}

\begin{itemize}
  \item \textbf{内环境}（细胞外液）：血浆、组织液、淋巴
  \item 内环境理化性质：渗透压（主要与 Na$^+$、Cl$^-$ 有关）、pH（7.35--7.45）、温度
  \item \textbf{稳态}：通过各器官系统协调，内环境保持相对稳定的状态
  \item 调节机制：神经--体液--免疫调节网络
\end{itemize}

\subsubsection{神经调节}

\begin{itemize}
  \item \textbf{反射}：神经调节的基本方式，反射弧（感受器 $\to$ 传入神经 $\to$ 神经中枢 $\to$ 传出神经 $\to$ 效应器）
  \item \textbf{兴奋传导}
  \begin{itemize}
    \item 静息电位：外正内负（K$^+$ 外流）
    \item 动作电位：外负内正（Na$^+$ 内流）
    \item 传导特点：双向传导（神经纤维上）
  \end{itemize}
  \item \textbf{突触传递}
  \begin{itemize}
    \item 结构：突触前膜、突触间隙、突触后膜
    \item 过程：电信号 $\to$ 化学信号（递质）$\to$ 电信号
    \item 传递特点：单向传递（突触延搁）
    \item 常见递质：乙酰胆碱（兴奋性）、多巴胺、$\gamma$-氨基丁酸（抑制性）
  \end{itemize}
  \item \textbf{分级调节}：大脑皮层（最高级中枢）$\to$ 脊髓（低级中枢）
  \item \textbf{语言中枢}：S 区（运动性失语症）、W 区（书写性失语症）、H 区（听觉性失语症）、V 区（视觉性失语症）
\end{itemize}

\subsubsection{体液调节}

\begin{itemize}
  \item \textbf{激素调节}
  \begin{itemize}
    \item 下丘脑：分泌促激素释放激素、抗利尿激素
    \item 垂体：分泌生长激素、促甲状腺激素、促性腺激素
    \item 甲状腺：甲状腺激素（促进新陈代谢、生长发育，提高神经兴奋性）
    \item 胰岛：A 细胞分泌胰高血糖素（升血糖），B 细胞分泌胰岛素（降血糖）
    \item 肾上腺：肾上腺素（使心跳加快、血压升高）
    \item 性腺：性激素（促进生殖器官发育、维持第二性征）
  \end{itemize}
  \item \textbf{反馈调节}
  \begin{itemize}
    \item 负反馈：维持稳态的主要方式（如甲状腺激素的分级调节）
    \item 正反馈：加速生理过程（如分娩、血液凝固）
  \end{itemize}
  \item 下丘脑--垂体--靶腺轴：分级调节
  \item 激素作用特点：微量高效、通过体液运输、作用于靶器官/靶细胞
\end{itemize}

\subsubsection{免疫调节}

\begin{itemize}
  \item \textbf{三道防线}
  \begin{itemize}
    \item 第一道：皮肤、黏膜
    \item 第二道：体液中的杀菌物质（溶菌酶）、吞噬细胞
    \item 第三道：特异性免疫（细胞免疫 + 体液免疫）
  \end{itemize}
  \item \textbf{体液免疫}：抗原 $\to$ B 细胞 $\xrightarrow{\text{辅助T细胞}}$ 浆细胞（抗体）+ 记忆 B 细胞
  \item \textbf{细胞免疫}：抗原 $\to$ 细胞毒性 T 细胞 $\xrightarrow{\text{辅助T细胞}}$ 细胞毒性 T 细胞（裂解靶细胞）+ 记忆 T 细胞
  \item \textbf{免疫异常}
  \begin{itemize}
    \item 过敏反应：已免疫机体再次接触相同过敏原
    \item 自身免疫病：免疫系统攻击自身组织（类风湿、系统性红斑狼疮）
    \item 免疫缺陷：HIV 攻击辅助 T 细胞 $\to$ 艾滋病
  \end{itemize}
  \item \textbf{疫苗}：灭活或减毒的病原体制成，刺激机体产生记忆细胞和抗体
\end{itemize}

\subsubsection{植物激素调节}

\begin{itemize}
  \item \textbf{生长素}
  \begin{itemize}
    \item 合成：幼芽、幼叶、发育中的种子（色氨酸）
    \item 运输：极性运输（形态学上端 $\to$ 下端，主动运输）
    \item 作用：两重性（低浓度促进生长，高浓度抑制生长）
    \item 顶端优势：顶芽优先生长，侧芽生长受抑制
  \end{itemize}
  \item \textbf{其他植物激素}
  \begin{itemize}
    \item 赤霉素：促进细胞伸长、种子萌发
    \item 细胞分裂素：促进细胞分裂
    \item 脱落酸：抑制细胞分裂、促进衰老脱落
    \item 乙烯：促进果实成熟
  \end{itemize}
  \item 植物激素间相互作用：协同/拮抗，共同调节生命活动
\end{itemize}

\subsection{生物与环境}

\subsubsection{种群}

\begin{itemize}
  \item 种群特征：出生率/死亡率、迁入率/迁出率、年龄结构、性别比例
  \item 种群密度调查方法：样方法（植物、活动能力弱动物）、标记重捕法（活动能力强动物）
  \item 种群增长模型
  \begin{itemize}
    \item J 型增长：理想条件，$N_t=N_0\lambda^t$
    \item S 型增长：有限资源，$K$ 为环境容纳量
  \end{itemize}
  \item 应用：$K/2$ 时增长速率最快（捕捞/狩猎最佳时期）
\end{itemize}

\subsubsection{群落}

\begin{itemize}
  \item 群落结构：垂直结构（分层）、水平结构（镶嵌）
  \item 种间关系
  \begin{itemize}
    \item 互利共生（+ +）：如根瘤菌与豆科植物
    \item 捕食（+ --）：如狼与羊
    \item 竞争（-- --）：如水稻与杂草
    \item 寄生（+ --）：如蛔虫与人
  \end{itemize}
  \item 群落演替
  \begin{itemize}
    \item 初生演替：裸岩 $\to$ 地衣 $\to$ 苔藓 $\to$ 草本 $\to$ 灌木 $\to$ 森林
    \item 次生演替：原有植被虽已不存在，但原有土壤条件基本保留
  \end{itemize}
\end{itemize}

\subsubsection{生态系统}

\begin{itemize}
  \item 组成成分：非生物的物质和能量、生产者（自养生物）、消费者（异养动物）、分解者（腐生细菌真菌）
  \item 营养结构：食物链（生产者为起点）和食物网
  \item \textbf{能量流动}
  \begin{itemize}
    \item 起点：生产者固定太阳能
    \item 特点：单向流动、逐级递减（传递效率 10--20\%）
    \item 去路：呼吸消耗、分解者利用、流入下一营养级、未利用
  \end{itemize}
  \item \textbf{物质循环}
  \begin{itemize}
    \item 碳循环：CO$_2$（无机环境）$\xrightarrow{\text{光合/化能合成}}$ 有机物（生物群落）$\xrightarrow{\text{呼吸/燃烧}}$ CO$_2$
    \item 温室效应：化石燃料大量燃烧，CO$_2$ 排放过多
  \end{itemize}
  \item \textbf{信息传递}
  \begin{itemize}
    \item 物理信息：光、声、温度、磁力
    \item 化学信息：信息素、外激素
    \item 行为信息：舞蹈、鸣叫
    \item 作用：利于生命活动正常进行、种群繁衍、调节种间关系
  \end{itemize}
  \item \textbf{生态系统稳定性}
  \begin{itemize}
    \item 抵抗力稳定性：抵抗干扰保持原状（与物种丰富度正相关）
    \item 恢复力稳定性：受到破坏后恢复原状
  \end{itemize}
\end{itemize}

\subsection{生物技术与工程}

\subsubsection{传统发酵技术}

\begin{itemize}
  \item 果酒制作：酵母菌（兼性厌氧），20℃ 左右，先通气后密封
  \item 果醋制作：醋酸菌（好氧），30--35℃，持续通气
  \item 泡菜制作：乳酸菌（厌氧），产生乳酸
  \item 腐乳制作：毛霉等微生物，加盐（抑制杂菌）+ 加酒（调味抑菌）
\end{itemize}

\subsubsection{微生物培养}

\begin{itemize}
  \item 培养基成分：水、碳源、氮源、无机盐（特殊微生物还需生长因子）
  \item 灭菌方法：灼烧灭菌（接种环）、干热灭菌（玻璃器皿）、湿热灭菌（培养基，高压蒸汽灭菌）
  \item 接种方法：平板划线法、稀释涂布平板法
  \item 计数方法：显微镜直接计数、稀释涂布平板法
\end{itemize}

\subsubsection{发酵工程}

\begin{itemize}
  \item 菌种选育 $\to$ 扩大培养 $\to$ 培养基配制灭菌 $\to$ 接种发酵 $\to$ 产物分离提纯
  \item 应用：柠檬酸、谷氨酸钠、青霉素、疫苗
\end{itemize}

\subsubsection{基因工程}

\begin{itemize}
  \item 操作工具
  \begin{itemize}
    \item 限制酶：切割 DNA 形成黏性末端或平末端
    \item DNA 连接酶：连接磷酸二酯键
    \item 载体：质粒、噬菌体、动植物病毒（需有复制原点、标记基因、多克隆位点）
  \end{itemize}
  \item 操作流程：目的基因获取 $\to$ 基因表达载体构建 $\to$ 导入受体细胞 $\to$ 筛选检测 $\to$ 目的基因表达
  \item 受体细胞：大肠杆菌（原核）、酵母菌、动植物细胞
  \item PCR 技术：变性（94℃）$\to$ 退火（50--60℃）$\to$ 延伸（72℃），需引物、Taq 酶、dNTP
  \item 应用：转基因作物、基因治疗、分子育种
\end{itemize}

\subsubsection{细胞工程}

\begin{itemize}
  \item \textbf{植物细胞工程}
  \begin{itemize}
    \item 植物组织培养：外植体 $\to$ 愈伤组织 $\to$ 再分化 $\to$ 完整植株，体现植物细胞全能性
    \item 植物体细胞杂交：去壁（纤维素酶、果胶酶）$\to$ 原生质体融合 $\to$ 杂种植株
    \item 应用：快速繁殖、脱毒苗、人工种子、突变体筛选
  \end{itemize}
  \item \textbf{动物细胞工程}
  \begin{itemize}
    \item 动物细胞培养：需动物血清、CO$_2$ 培养箱（5\% CO$_2$ 维持 pH）
    \item 核移植技术：体细胞核 + 去核卵母细胞 $\to$ 重组细胞 $\to$ 克隆动物
    \item 单克隆抗体：B 细胞 + 骨髓瘤细胞 $\to$ 杂交瘤细胞 $\to$ 单克隆抗体
    \item 动物细胞融合：PEG、灭活病毒、电融合
  \end{itemize}
\end{itemize}

\subsubsection{胚胎工程}

\begin{itemize}
  \item 体外受精：卵母细胞培养 + 精子获能 $\to$ 受精
  \item 胚胎移植：供体超数排卵（促性腺激素）$\to$ 冲卵 $\to$ 移植到受体
  \item 胚胎分割：桑椹胚或囊胚期进行，注意内细胞团均等分割
  \item 应用：转基因动物、试管动物、胚胎干细胞研究
\end{itemize}

\subsection{实验方法}

\subsubsection{显微镜观察类实验}

\begin{itemize}
  \item \textbf{观察 DNA 和 RNA 在细胞中的分布}：甲基绿（DNA 绿）、吡罗红（RNA 红），8\% 盐酸改变膜通透性、使 DNA 与蛋白质分离
  \item \textbf{观察线粒体}：健那绿染液（活体染色，线粒体蓝绿色）
  \item \textbf{观察叶绿体}：清水制片，不需染色
  \item \textbf{观察质壁分离与复原}：紫色洋葱鳞片叶外表皮，0.3 g/mL 蔗糖溶液
  \item \textbf{观察细胞有丝分裂}：洋葱根尖分生区，解离$\to$漂洗$\to$染色$\to$制片（龙胆紫/醋酸洋红）
  \item \textbf{观察减数分裂}：蝗虫精巢固定装片
\end{itemize}

\subsubsection{生化提取与鉴定类实验}

\begin{itemize}
  \item \textbf{还原糖鉴定}：斐林试剂（NaOH + CuSO$_4$），水浴加热，砖红色沉淀
  \item \textbf{脂肪鉴定}：苏丹 III（橘黄色）/ 苏丹 IV（红色）
  \item \textbf{蛋白质鉴定}：双缩脲试剂（先加 A 液 NaOH，后加 B 液 CuSO$_4$），紫色反应
  \item \textbf{淀粉鉴定}：碘液，蓝色
  \item \textbf{DNA 鉴定}：二苯胺试剂，沸水浴，蓝色
  \item \textbf{绿叶中色素的提取和分离}
  \begin{itemize}
    \item 提取：无水乙醇（溶剂），SiO$_2$（研磨充分），CaCO$_3$（防止色素被破坏）
    \item 分离：纸层析法，层析液为石油醚+丙酮+苯
  \end{itemize}
\end{itemize}

\subsubsection{探究类实验}

\begin{itemize}
  \item \textbf{探究酶活性影响因素}：控制单一变量，设置温度/pH 梯度
  \item \textbf{探究酵母菌呼吸方式}：澄清石灰水/溴麝香草酚蓝（CO$_2$）、重铬酸钾（酒精，灰绿色）
  \item \textbf{探究生长素类似物促进插条生根最适浓度}：预实验 $\to$ 细化浓度梯度
  \item \textbf{探究培养液中酵母菌种群数量变化}：血细胞计数板计数，每天定时取样
  \item \textbf{探究土壤微生物分解作用}：对照实验，控制变量
\end{itemize}

\subsubsection{实验设计基本原则}

\begin{itemize}
  \item \textbf{对照原则}：空白对照、自身对照、相互对照、条件对照
  \item \textbf{单一变量原则}：只有一个自变量，其他条件相同且适宜
  \item \textbf{平行重复原则}：每组至少 3 个重复，减少偶然误差
  \item \textbf{随机性原则}：样本选取随机
\end{itemize}

\subsubsection{常用试剂与作用速查}

\[
\begin{array}{c|c}
\text{试剂} & \text{作用} \\ \hline
斐林试剂 & 鉴定还原糖（砖红色沉淀） \\
双缩脲试剂 & 鉴定蛋白质（紫色） \\
苏丹 III/IV & 脂肪染色（橘黄/红色） \\
龙胆紫/醋酸洋红 & 染色体染色（紫色/红色） \\
健那绿 & 线粒体活体染色（蓝绿色） \\
甲基绿/吡罗红 & DNA（绿）/RNA（红） \\
无水乙醇 & 提取色素 \\
层析液 & 分离色素 \\
秋水仙素 & 抑制纺锤体形成（诱导多倍体） \\
碳酸钙 & 保护色素 \\
二苯胺试剂 & DNA 鉴定（蓝色） \\
重铬酸钾 & 酒精检测（灰绿色）
\end{array}
\]
